\documentclass[aspectratio=169]{beamer}
\usetheme{Madrid}
\usecolortheme{default}
\usepackage{graphicx}
\usepackage{amsmath}
\usepackage{booktabs}
\usepackage{hyperref}
% UGA Colors
\definecolor{UGARed}{HTML}{BA0C2F}
\definecolor{UGABlack}{HTML}{000000}
% Apply UGA colors to Madrid theme
\setbeamercolor{palette primary}{bg=UGARed,fg=white}
\setbeamercolor{palette secondary}{bg=UGABlack,fg=white}
\setbeamercolor{palette tertiary}{bg=UGARed,fg=white}
\setbeamercolor{title}{fg=UGARed}
\setbeamercolor{frametitle}{bg=UGARed,fg=white}
\setbeamercolor{structure}{fg=UGARed}
% UGA Logo on title page
\titlegraphic{\includegraphics[height=1.5cm]{download.png}}
\logo{\includegraphics[height=0.5cm]{download.png}}
% Title page information
\title{Privacy, Surveillance, and Ethics}
\subtitle{From Government Tool to Pervasive Force}
\author{Breakout Discussion Session}
\date{February 2026}
\begin{document}
% Slide 1: Title
\begin{frame}
    \titlepage
\end{frame}
% Slide 2: Overview / Context
\begin{frame}{This Week's Theme}
    \large
    \textbf{A common argument:} ``If you have nothing to hide, you have nothing to fear.''
    \vspace{0.4cm}
    Surveillance began as a mundane bureaucratic tool: counting people,
    planning cities, allocating budgets, and managing transport.
    \vspace{0.3cm}
    Ethically, the ``nothing to hide'' view is widely considered
    a \textbf{false trade-off}. Privacy protects ordinary,
    lawful behavior and individual dignity, not just wrongdoing.
    \vspace{0.3cm}
    \textbf{This week's discussion} traces how those tools for
    organizing populations evolved into pervasive digital surveillance.
    \vspace{0.2cm}
    \begin{center}
        \href{https://youtu.be/I1NJyVNERBw?si=g2IoqHjLhF2fcxMy}{\textcolor{UGARed}{Listen: Surveillance Overview (Podcast)}}
    \end{center}
    \vspace{0.2cm}
    \textbf{\textcolor{UGARed}{Central question:}} Is surveillance still the same tool it once was?
\end{frame}
% Slide 3: Question 1
\begin{frame}{Privacy vs. Security: A Real Trade-Off?}
    \normalsize
    \textit{Surveillance started as a tool to count people, plan cities, and manage public health.
    After 9/11 it expanded rapidly into policing, private companies, and everyday life.}
    \vspace{0.3cm}
    \begin{itemize}
        \setlength\itemsep{0.3cm}
        \item 9/11 opened the floodgates for surveillance framed as security. Was that trade-off ever really a choice for ordinary people?
        \item Private companies sometimes hold more data than governments, tracking your phone, your searches, your location. Is that still about security, or is it about profit?
    \end{itemize}
\end{frame}
% Slide 4: Question 2
\begin{frame}{Should Platforms Explain How Data is Used?}
    \normalsize
    \textit{Clicks, likes, location, and shares are collected, not just to improve the service,
    but to predict behavior and sell those predictions.}
    \vspace{0.3cm}
    \begin{itemize}
        \setlength\itemsep{0.3cm}
        \item ``We don't sell your data, we sell the insights from it.'' What does that actually mean, and do users understand that distinction?
        \item Google used search data to improve search, then realized the leftover digital exhaust was worth more as a product sold to advertisers. Were users ever told that?
    \end{itemize}
\end{frame}
% Slide 5: Question 3
\begin{frame}{Who is Responsible for Protecting Personal Data?}
    \normalsize
    \textit{Governments outsource surveillance to private firms. Private firms profit and expand.}
    \vspace{0.3cm}
    \begin{itemize}
        \setlength\itemsep{0.3cm}
        \item When responsibility is shared between governments and corporations like that, who actually owns it?
        \item Can individuals realistically protect themselves from tuning, herding, and conditioning: subtle behavioral manipulation they cannot see?
    \end{itemize}
\end{frame}
% Slide 6: Question 4
\begin{frame}{How Do We Fix the Power Imbalance?}
    \normalsize
    \textit{Doctors are expected to use patient data only for treatment.
    Supermarkets should not sell purchase habits to insurance companies.
    That line is already being crossed, quietly and legally.}
    \vspace{0.3cm}
    \begin{itemize}
        \setlength\itemsep{0.3cm}
        \item When systems classify people into risk groups, who gets watched more, and does the data make that fair or just make it feel fair?
        \item Facial recognition wrongly identifies African-Americans at significantly higher rates in serious criminal cases. Who bears the cost of that imbalance?
    \end{itemize}
\end{frame}
% Slide 7: Question 5
\begin{frame}{When Can Privacy Be Justifiably Limited?}
    \normalsize
    \textit{COVID tracing apps were built for public health.
    Malcolm Badali blogged about migrant workers in Qatar and was arrested.}
    \vspace{0.3cm}
    \begin{itemize}
        \setlength\itemsep{0.3cm}
        \item At what point does a legitimate safety measure become a surveillance tool?
        \item Who should have the authority to draw that line, and how do we hold them to it?
    \end{itemize}
\end{frame}
% Slide 8: Closing
\begin{frame}{Closing Thought}
    \large
    \vspace{0.5cm}
    \centering
    In a world where our experiences are treated as raw material\\
    and privacy feels like it is shrinking...
    \vspace{0.6cm}
    \textbf{\textcolor{UGARed}{What does it really mean to have a digital home?}}
    \vspace{0.4cm}
    \textbf{What can we do, individually and together,\\
    to reclaim some sense of safety in our digital lives?}
    \vspace{0.6cm}
    \normalsize
    \textit{``Privacy is not just about secrets. It is about power, dignity,\\
    and the freedom to live without being watched.''}\\
    \vspace{0.2cm}
    \textit{Daniel J. Solove}
\end{frame}
\end{document}
% ============================================================================
% PRESENTATION STRUCTURE (8 slides total):
% ============================================================================
%
% Slide 1:  Title
% Slide 2:  Theme overview / context + podcast link
% Slide 3:  Q1 -- Nothing to hide / security trade-off
% Slide 4:  Q2 -- Platform transparency / behavioral surplus
% Slide 5:  Q3 -- Responsibility: individuals, corporations, governments
% Slide 6:  Q4 -- Power imbalances / contextual integrity
% Slide 7:  Q5 -- Justified limits / Badali / COVID apps
% Slide 8:  Closing -- Digital home + Solove
%
% ============================================================================
